% ============================================================
%  DISPENSA UNIVERSITARIA
%  Funzionamento dei Sistemi Biologici
% ============================================================
\documentclass[11pt,a4paper,openright,twoside]{book}

% --- Encoding & Language ---
\usepackage[utf8]{inputenc}
\usepackage[T1]{fontenc}
\usepackage[italian]{babel}

% --- Fonts ---
\usepackage{lmodern}
\usepackage{microtype}

% --- Math ---
\usepackage{amsmath,amssymb,amsthm}
\usepackage{mathtools}

% --- Layout ---
\usepackage[
  a4paper,
  top=1.5cm, bottom=1.5cm,
  inner=1.5cm, outer=1.5cm,
  headsep=0.8cm
]{geometry}

% --- Headers & Footers ---
\usepackage{fancyhdr}
\pagestyle{fancy}
\fancyhf{}
\fancyhead[LE]{\small\nouppercase{\leftmark}}
\fancyhead[RO]{\small\nouppercase{\rightmark}}
\fancyfoot[LE,RO]{\thepage}
\renewcommand{\headrulewidth}{0.4pt}

% --- Colors & Boxes ---
\usepackage[dvipsnames,svgnames]{xcolor}
\usepackage[most]{tcolorbox}
\tcbuselibrary{breakable,skins}

% Define color palette
\definecolor{PrimaryBlue}{RGB}{0,74,127}
\definecolor{AccentTeal}{RGB}{0,128,128}
\definecolor{LightGray}{RGB}{245,245,250}
\definecolor{DefinitionBg}{RGB}{230,244,255}
\definecolor{ExampleBg}{RGB}{240,255,240}
\definecolor{NoteBg}{RGB}{255,250,230}
\definecolor{WarningBg}{RGB}{255,235,235}

% --- Tables ---
\usepackage{booktabs}
\usepackage{array}
\usepackage{tabularx}
\usepackage{multirow}

% --- Graphics ---
\usepackage{graphicx}
\usepackage{tikz}
\usetikzlibrary{arrows.meta,shapes,positioning,decorations.pathmorphing,calc}
\usepackage{caption}
\usepackage{subcaption}

% --- Lists ---
\usepackage{enumitem}
\setlist[itemize]{itemsep=2pt, topsep=4pt}
\setlist[enumerate]{itemsep=2pt, topsep=4pt}

% --- Section Titles ---
\usepackage{titlesec}
\titleformat{\chapter}[display]
  {\normalfont\huge\bfseries\color{PrimaryBlue}}
  {\chaptername\ \thechapter}{18pt}{\Huge}
\titleformat{\section}
  {\normalfont\Large\bfseries\color{PrimaryBlue}}
  {\thesection}{1em}{}
\titleformat{\subsection}
  {\normalfont\large\bfseries\color{AccentTeal}}
  {\thesubsection}{1em}{}
\titleformat{\subsubsection}
  {\normalfont\normalsize\bfseries}
  {\thesubsubsection}{1em}{}

% --- Hyperref ---
\usepackage[
  colorlinks=true,
  linkcolor=PrimaryBlue,
  citecolor=AccentTeal,
  urlcolor=AccentTeal,
  bookmarksnumbered=true,
  pdftitle={Funzionamento dei sistemi biologici}
]{hyperref}
\usepackage{cleveref}

% --- Theorem Environments ---
\theoremstyle{definition}
\newtheorem{definizione}{Definizione}[chapter]
\newtheorem{esempio}{Esempio}[chapter]
\theoremstyle{plain}
\newtheorem{proprieta}{Proprietà}[chapter]
\theoremstyle{remark}
\newtheorem{osservazione}{Osservazione}[chapter]
\newtheorem{nota}{Nota}[chapter]

% --- Custom tcolorbox styles ---
\tcbset{
  defbox/.style={
    enhanced, breakable,
    colback=DefinitionBg, colframe=PrimaryBlue,
    fonttitle=\bfseries, title={Definizione},
    left=6pt, right=6pt, top=4pt, bottom=4pt,
    arc=4pt
  },
  exbox/.style={
    enhanced, breakable,
    colback=ExampleBg, colframe=AccentTeal,
    fonttitle=\bfseries, title={Esempio},
    left=6pt, right=6pt, top=4pt, bottom=4pt,
    arc=4pt
  },
  notebox/.style={
    enhanced, breakable,
    colback=NoteBg, colframe=orange!70!black,
    fonttitle=\bfseries, title={Nota},
    left=6pt, right=6pt, top=4pt, bottom=4pt,
    arc=4pt
  },
  summarybox/.style={
    enhanced, breakable,
    colback=LightGray, colframe=PrimaryBlue!60,
    fonttitle=\bfseries\color{PrimaryBlue},
    left=8pt, right=8pt, top=6pt, bottom=6pt,
    arc=6pt
  }
}

% --- Utility commands ---
\newcommand{\dna}{\textsc{dna}}
\newcommand{\rna}{\textsc{rna}}
\newcommand{\mrna}{m\textsc{rna}}
\newcommand{\trna}{t\textsc{rna}}
\newcommand{\rrna}{r\textsc{rna}}
\newcommand{\atp}{\textsc{atp}}
\newcommand{\adp}{\textsc{adp}}
\newcommand{\nadh}{\textsc{nadh}}
\newcommand{\nadph}{\textsc{nadph}}
\newcommand{\pH}{\ensuremath{\text{pH}}}
\newcommand{\Ca}{\ensuremath{\text{Ca}^{2+}}}
\newcommand{\Na}{\ensuremath{\text{Na}^{+}}}
\newcommand{\K}{\ensuremath{\text{K}^{+}}}
\newcommand{\CO}{\ensuremath{\text{CO}_2}}
\newcommand{\HO}{\ensuremath{\text{H}_2\text{O}}}
\newcommand{\Oc}{\ensuremath{\text{O}_2}}

% ============================================================
\begin{document}
% ============================================================

% --- Title Page ---
\begin{titlepage}
  \centering
  \vspace*{2cm}
  {\color{PrimaryBlue}\rule{\linewidth}{1pt}}\\[0.5cm]
  {\Huge\bfseries\color{PrimaryBlue} Funzionamento dei sistemi biologici}\\[0.5cm]
  {\color{PrimaryBlue}\rule{\linewidth}{1pt}}\\[0.5cm]
  \vfill
\end{titlepage}

\frontmatter

% --- Table of Contents ---
\tableofcontents
\cleardoublepage

\mainmatter

% ============================================================
\chapter{Organismi Viventi}
% ============================================================

\begin{tcolorbox}[summarybox, title={Sommario del capitolo}]
Questo capitolo introduce i fondamenti della biologia: la classificazione degli organismi, le principali teorie evolutive, la teoria cellulare, la composizione chimica della materia vivente (carboidrati, lipidi, proteine, acidi nucleici) e le caratteristiche funzionali degli organismi. Vengono inoltre trattati la citologia (virus, batteriofagi, cellule procariotiche) e le principali biotecnologie.
\end{tcolorbox}

% ============================================================
\section{Biologia}
% ============================================================

Il termine \textbf{biologia} deriva dal greco (\textit{bìos}, vita; \textit{lògos}, studio) e designa la scienza che studia i fenomeni della vita in tutte le loro manifestazioni. La biologia si articola attorno a tre pilastri fondamentali:
\begin{enumerate}
  \item \textbf{Evoluzione:} tutte le specie viventi derivano da forme primordiali attraverso processi evolutivi graduali;
  \item \textbf{Trasmissione dell'informazione genetica:} ogni organismo trasmette alla propria discendenza le istruzioni necessarie allo sviluppo e al funzionamento;
  \item \textbf{Trasferimento di energia:} i processi biologici richiedono e producono energia, prevalentemente veicolata dalla molecola \atp{} (adenosin trifosfato).
\end{enumerate}

\subsection{Livelli di organizzazione della materia vivente}

L'organizzazione biologica procede per livelli gerarchici progressivamente più complessi, a partire dagli atomi:

\[
\text{Atomi} \;\to\; \text{Molecole} \;\to\; \text{Macromolecole} \;\to\; \text{Cellule} \;\to\; \text{Tessuti} \;\to\; \text{Organi} \;\to\; \text{Apparati}
\]
\[
\to\; \text{Organismi} \;\to\; \text{Popolazioni} \;\to\; \text{Comunità} \;\to\; \text{Ecosistemi} \;\to\; \text{Biosfera}
\]

Le principali \textbf{macromolecole biologiche} sono:
\begin{itemize}
  \item Carboidrati (o glicidi)
  \item Lipidi
  \item Proteine
  \item Acidi nucleici
\end{itemize}

Le cellule si distinguono in due grandi categorie:
\begin{itemize}
  \item \textbf{Cellule procarioti:} prive di nucleo definito; il materiale genetico è disperso nel citoplasma; contengono ribosomi racchiusi dalla membrana plasmatica;
  \item \textbf{Cellule eucarioti:} dotate di nucleo delimitato da involucro nucleare; il materiale genetico è organizzato in cromosomi; si specializzano in funzioni specifiche formando i tessuti degli organismi pluricellulari.
\end{itemize}

% ============================================================
\subsection{Classificazione degli organismi}
% ============================================================

La tassonomia classifica gli organismi mediante una gerarchia di livelli sistematici, dal più generale al più specifico:

\[
\text{Dominio} \to \text{Regno} \to \text{Phylum} \to \text{Classe} \to \text{Ordine} \to \text{Famiglia} \to \text{Genere} \to \text{Specie}
\]

\begin{tcolorbox}[exbox, title={Esempio: classificazione dell'essere umano}]
  \textit{Homo sapiens} appartiene al dominio \textit{Eukarya}, regno \textit{Animalia}, phylum \textit{Chordata}, classe \textit{Mammalia}, ordine \textit{Primates}, famiglia \textit{Hominidae}, genere \textit{Homo}, specie \textit{sapiens}.
\end{tcolorbox}

% ============================================================
\subsection{Teoria dell'evoluzione}
% ============================================================

Nel corso della storia della biologia sono state elaborate diverse teorie per spiegare l'origine e la diversità delle specie.

\subsubsection{Teoria di Linneo}
Linneo sosteneva l'esistenza di specie fisse e immutabili nel tempo. Questa concezione fu gradualmente confutata grazie al ritrovamento di fossili di organismi non più presenti sulla Terra, a dimostrazione che le specie possono estinguersi e che la vita si è modificata nel tempo.

\subsubsection{Teoria di Lamarck}
Lamarck propose che gli organismi si modifichino in risposta alle condizioni ambientali e che tali modificazioni vengano trasmesse alla prole (\textit{ereditarietà dei caratteri acquisiti}). Questa teoria fu successivamente confutata, poiché le modificazioni fenotipiche acquisite durante la vita di un individuo non sono ereditabili.

\subsubsection{Teoria di Darwin}
Charles Darwin elaborò la teoria dell'evoluzione per selezione naturale, basata su tre osservazioni fondamentali:
\begin{enumerate}
  \item \textbf{Variabilità:} gli individui di una popolazione mostrano variazioni nelle loro caratteristiche;
  \item \textbf{Selezione naturale:} nelle condizioni di competizione per risorse limitate, gli individui meglio adattati all'ambiente hanno maggiori probabilità di sopravvivere e riprodursi;
  \item \textbf{Ereditarietà:} le caratteristiche favorevoli vengono trasmesse alla prole attraverso una base genetica.
\end{enumerate}

Darwin osservò, inoltre, che la maggior parte degli organismi produce più discendenti di quanti ne sopravvivano, che le popolazioni tendono a rimanere stabili nel tempo e che le risorse disponibili nell'ambiente sono limitate. Da queste premesse dedusse che gli individui competono per tali risorse e che la selezione agisce su riproduzione e sopravvivenza. Darwin elaborò anche la teoria dell'\textit{albero della vita}, secondo cui tutti i viventi condividono un antenato primordiale comune.

% ============================================================
\subsection{Teoria cellulare}
% ============================================================

\begin{tcolorbox}[defbox, title={Teoria cellulare}]
  La teoria cellulare afferma che:
  \begin{enumerate}
    \item tutti gli organismi viventi sono composti da una o più cellule;
    \item la cellula è l'unità strutturale e funzionale fondamentale della vita;
    \item ogni cellula deriva da una cellula preesistente per divisione cellulare.
  \end{enumerate}
\end{tcolorbox}

Ogni cellula è caratterizzata da quattro proprietà fondamentali:
\begin{enumerate}
  \item \textbf{Metabolismo:} insieme delle reazioni chimiche che avvengono al suo interno, necessarie al mantenimento della vita;
  \item \textbf{Accrescimento:} capacità di assimilare nutrienti e aumentare di dimensione;
  \item \textbf{Riproduzione:} capacità di duplicarsi, generando cellule figlie;
  \item \textbf{Irritabilità:} capacità di percepire e rispondere a stimoli esterni (temperatura, luce, stimoli chimici).
\end{enumerate}

Dal punto di vista chimico, la cellula è composta per circa il 75--85\% di acqua e per il 15--25\% da componenti organiche (carboidrati, lipidi, proteine, acidi nucleici) e inorganiche (elementi primari C, H, O, N, S, P; cationi e anioni; oligoelementi).

\subsection{Ordini di grandezza e strumenti di osservazione}

Per osservare strutture biologiche a diverse scale di ingrandimento si utilizzano diversi strumenti, ciascuno con un proprio \textit{limite di risoluzione} (distanza minima al di sotto della quale due punti distinti non possono essere discriminati):

\begin{center}
\begin{tabular}{lll}
\toprule
\textbf{Strumento} & \textbf{Limite di risoluzione} & \textbf{Applicazioni tipiche} \\
\midrule
Occhio umano & $\sim 0{,}1\;\text{mm}$ & Osservazione macroscopica \\
Microscopio ottico & $\sim 0{,}2\;\mu\text{m}$ & Osservazione di cellule \\
Microscopio elettronico & $\sim 0{,}1\;\text{nm}$ & Osservazione di organuli \\
\bottomrule
\end{tabular}
\end{center}

% ============================================================
\section{La materia vivente: le macromolecole biologiche}
% ============================================================

\subsection{Carboidrati}

I \textbf{carboidrati} (detti anche glicidi o zuccheri) rappresentano la principale fonte di energia per le cellule, contribuendo per oltre la metà dell'apporto energetico giornaliero (250--800\,g al giorno nell'essere umano). Sono costituiti da carbonio, idrogeno e ossigeno, in rapporto generale $\text{C}_n(\text{H}_2\text{O})_n$.

\subsubsection{Classificazione}

L'unità fondamentale dei carboidrati è il \textbf{monosaccaride}, classificato in base al numero di atomi di carbonio: triosi (3C), pentosi (5C), esosi (6C).

La condensazione di più monosaccaridi forma gli \textbf{oligosaccaridi}. Tra questi, i \textbf{disaccaridi} sono formati dall'unione di due monosaccaridi:
\begin{itemize}
  \item Saccarosio = glucosio + fruttosio
  \item Maltosio = glucosio + glucosio
  \item Lattosio = glucosio + galattosio
\end{itemize}

I \textbf{polisaccaridi} sono polimeri di monosaccaridi e si dividono in:
\begin{itemize}
  \item \textit{Omopolisaccaridi:} composti dalla stessa unità monosaccaridica;
  \item \textit{Eteropolisaccaridi:} composti da unità monosaccaridiche diverse.
\end{itemize}

\subsubsection{Principali polisaccaridi}

\textbf{Amido e glicogeno} sono i principali polisaccaridi di deposito, rispettivamente nelle cellule vegetali e animali. Entrambi sono formati da:
\begin{itemize}
  \item \textit{Amilosio:} catena lineare di molecole di glucosio unite da legami $\alpha[1\!\to\!4]$ monoglicosilici;
  \item \textit{Amilopectina:} catena ramificata con legami $\alpha[1\!\to\!6]$ glicosilici nei punti di ramificazione.
\end{itemize}

\textbf{Cellulosa:} polimero lineare di migliaia di monomeri di glucosio uniti da legami $\beta[1\!\to\!4]$. Le catene si dispongono parallelamente tramite legami a idrogeno, formando \textit{fibrille}. Ha funzione strutturale nei vegetali, costituendo la parete cellulare primaria e secondaria. Nell'organismo umano non è digeribile dagli enzimi, ma viene fermentata dalla flora batterica intestinale.

% ============================================================
\subsection{Lipidi}
% ============================================================

I \textbf{lipidi} (comunemente detti grassi) svolgono funzioni di riserva energetica, strutturale e di segnalazione. Nell'organismo si suddividono in:
\begin{itemize}
  \item \textbf{Trigliceridi} (98\% dei lipidi totali): lipidi di deposito con funzione di riserva energetica, costituiscono il tessuto adiposo;
  \item \textbf{Fosfolipidi, glicolipidi, colesterolo e vitamine liposolubili} (2\%): svolgono funzione strutturale nelle membrane cellulari; dal colesterolo si sintetizzano sali biliari e vitamina D.
\end{itemize}

\subsubsection{Acidi grassi}

Sono catene lineari con numero pari di atomi di carbonio e si classificano in:
\begin{enumerate}
  \item \textbf{Acidi grassi saturi:} legami semplici tra tutti gli atomi di carbonio; la solubilità in acqua diminuisce all'aumentare della lunghezza della catena;
  \item \textbf{Acidi grassi insaturi:} presentano uno o più doppi legami nella catena:
    \begin{itemize}
      \item \textit{Monoinsaturi:} un solo doppio legame (es. acido oleico);
      \item \textit{Polinsaturi:} due o più doppi legami (es. acido linoleico); la solubilità in acqua è direttamente proporzionale al numero di doppi legami.
    \end{itemize}
\end{enumerate}

\subsubsection{Lipidi semplici e complessi}

I \textbf{lipidi semplici} derivano dall'esterificazione di acidi grassi con alcoli. I \textbf{trigliceridi}, ad esempio, derivano dall'esterificazione del glicerolo con 3 acidi grassi (con eliminazione di 3 molecole d'acqua) e svolgono funzione di riserva energetica, isolamento termico e sostegno degli organi.

I \textbf{lipidi complessi} derivano dai trigliceridi per sostituzione di un acido grasso con un gruppo fosfato e altri residui:
\begin{itemize}
  \item \textbf{Fosfolipidi:} testa idrofila (gruppo fosfato) e code idrofobiche (acidi grassi); costituiscono il doppio strato delle membrane biologiche;
  \item \textbf{Steroidi:} struttura a 4 anelli fusi con un gruppo ossidrilico; comprendono colesterolo e vitamina D.
\end{itemize}

% ============================================================
\subsection{Proteine}
% ============================================================

Le \textbf{proteine} sono le biomolecole più versatili degli organismi viventi. Sono polimeri di \textbf{amminoacidi}, molecole bipolari formate da un gruppo amminico ($-\text{NH}_2$), un gruppo carbossilico ($-\text{COOH}$), un atomo di idrogeno e una catena laterale variabile (gruppo R).

Gli amminoacidi proteinogenici sono \textbf{20}, classificabili in base alla carica:
\begin{itemize}
  \item Non polari (idrofobici);
  \item Polari non carichi;
  \item Carichi negativamente (acidi);
  \item Carichi positivamente (basici).
\end{itemize}

In base alla necessità nutrizionale, si distinguono in:
\begin{itemize}
  \item \textbf{Essenziali:} non sintetizzabili dall'organismo in quantità sufficiente (devono essere introdotti con la dieta);
  \item \textbf{Condizionatamente essenziali:} sintetizzabili, ma non a sufficienza in particolari condizioni fisiologiche;
  \item \textbf{Non essenziali:} sintetizzabili autonomamente a partire da intermedi metabolici.
\end{itemize}

Gli amminoacidi si legano mediante \textbf{legami peptidici}, che si formano per condensazione tra il gruppo carbossilico di un amminoacido e il gruppo amminico del successivo, con eliminazione di una molecola d'acqua.

\subsubsection{Livelli di struttura proteica}

\begin{enumerate}
  \item \textbf{Struttura primaria:} sequenza lineare degli amminoacidi nella catena polipeptidica;
  \item \textbf{Struttura secondaria:} ripiegamento locale della catena stabilizzato da legami a idrogeno tra il backbone peptidico:
    \begin{itemize}
      \item \textit{$\alpha$-elica:} la catena si avvolge in spirale con i gruppi R proiettati all'esterno;
      \item \textit{Foglietto $\beta$:} la catena si piega su sé stessa, con i gruppi R disposti alternamente sopra e sotto il piano;
    \end{itemize}
  \item \textbf{Struttura terziaria:} conformazione tridimensionale globale, stabilizzata dalle interazioni tra i gruppi R (legami a idrogeno, ionici, interazioni idrofobiche, ponti disolfuro);
  \item \textbf{Struttura quaternaria:} assemblaggio di più catene polipeptidiche in un complesso funzionale; esempi:
    \begin{itemize}
      \item \textit{Emoglobina:} 4 catene (2 $\alpha$-eliche e 2 $\beta$-foglietti);
      \item \textit{Collagene:} 3 catene identiche a formare una tripla elica.
    \end{itemize}
\end{enumerate}

\subsubsection{Funzioni delle proteine}

\begin{center}
\begin{tabular}{ll}
\toprule
\textbf{Funzione} & \textbf{Esempi} \\
\midrule
Strutturale & Cheratina, collagene \\
Enzimatica & Amilasi, DNA polimerasi \\
Trasporto & Emoglobina, albumina \\
Motrice & Actina, miosina \\
Regolatrice & Istoni, fattori di trascrizione \\
Recettoriale & Recettori di membrana \\
Ormonale & Insulina, ormone della crescita \\
Difesa & Anticorpi (immunoglobuline) \\
Accumulo & Ferritina \\
\bottomrule
\end{tabular}
\end{center}

% ============================================================
\subsection{Acidi nucleici}
% ============================================================

Gli \textbf{acidi nucleici} sono polimeri di nucleotidi, ognuno dei quali è composto da:
\begin{itemize}
  \item Uno zucchero pentoso;
  \item Una base azotata;
  \item Uno o più gruppi fosfato.
\end{itemize}

Le basi azotate si suddividono in:
\begin{itemize}
  \item \textbf{Pirimidine} (anello semplice): citosina (DNA e \rna), timina (solo \dna), uracile (solo \rna);
  \item \textbf{Purine} (doppio anello biciclico): adenina e guanina (presenti in entrambi).
\end{itemize}

Le basi si appaiano secondo regole precise: adenina--timina (o uracile nell'\rna) mediante 2 legami a idrogeno; guanina--citosina mediante 3 legami a idrogeno.

\subsubsection{DNA (Acido Desossiribonucleico)}

Il \dna{} è formato da nucleotidi contenenti lo zucchero \textit{desossiribosio}, un gruppo fosfato e una delle quattro basi (adenina, timina, guanina, citosina). Due catene antiparallele si avvolgono a \textbf{doppia elica}, stabilizzata dai legami a idrogeno tra le basi complementari. Le funzioni del \dna{} sono:
\begin{itemize}
  \item Conservazione e trasmissione dei caratteri ereditari;
  \item Guida per la propria duplicazione;
  \item Controllo della sintesi proteica.
\end{itemize}

Nelle cellule, la catena di \dna{} (molto lunga) è condensata mediante super-avvolgimento grazie agli \textbf{istoni}, proteine basiche (cariche positivamente) che si associano a gruppi di 8 formando il nucleosoma, interagendo con il \dna{} carico negativamente.

La \textbf{replicazione del \dna} è un processo \textit{semiconservativo}: le due catene si separano e la DNA-polimerasi sintetizza una catena complementare per ciascuna, ottenendo così due molecole figlie, ognuna con un filamento originale e uno neo-sintetizzato.

\subsubsection{RNA (Acido Ribonucleico)}

L'\rna{} è formato da nucleotidi contenenti il ribosio, un gruppo fosfato e una delle quattro basi (adenina, uracile, guanina, citosina). Esistono tre principali tipi funzionali:

\begin{itemize}
  \item \textbf{mRNA} (RNA messaggero): trasporta l'informazione codificata nel \dna{} ai ribosomi, dove viene tradotta in sequenza amminoacidica;
  \item \textbf{tRNA} (RNA di trasporto): trasporta gli amminoacidi al ribosoma, riconoscendo i codoni dell'mRNA tramite l'anticodone;
  \item \textbf{rRNA} (RNA ribosomiale): componente strutturale e catalitica dei ribosomi, essenziale per la decodifica dell'mRNA e la formazione dei legami peptidici.
\end{itemize}

% ============================================================
\section{Gli organismi viventi: caratteristiche generali}
% ============================================================

Un organismo vivente è definito dalla capacità di:
\begin{enumerate}
  \item Riprodursi e trasmettere i propri caratteri ereditari;
  \item Compiere un ciclo vitale (nascita, crescita, sviluppo, morte);
  \item Rinnovare continuamente la propria struttura tramite il metabolismo;
  \item Reagire agli stimoli ambientali;
  \item Muoversi (attivamente o passivamente);
  \item Evolversi nel tempo geologico.
\end{enumerate}

\subsection{Riproduzione}

La \textbf{riproduzione} è l'insieme dei meccanismi che assicurano la conservazione della specie e si divide in:

\begin{itemize}
  \item \textbf{Asessuata:} non comporta modificazione del patrimonio ereditario; il genotipo della progenie è identico a quello del genitore. Include:
    \begin{itemize}
      \item \textit{Moltiplicazione cellulare} (divisione mitotica): un organismo unicellulare si divide in due cellule figlie geneticamente identiche;
      \item \textit{Scissione:} tipica dei procarioti;
      \item \textit{Frammentazione:} parti staccate del corpo ricostituiscono un individuo intero;
      \item \textit{Gemmazione:} formazione di gemme laterali che originano nuovi individui;
      \item \textit{Poliembrionia:} divisione dello zigote nelle prime fasi di sviluppo (gemelli omozigoti);
      \item \textit{Sporulazione:} produzione di spore resistenti in condizioni ambientali avverse.
    \end{itemize}
  \item \textbf{Sessuata:} prevede la fusione (fecondazione) di due gameti (uno maschile, uno femminile), prodotti per meiosi, generando variabilità genetica. Si distingue in:
    \begin{itemize}
      \item \textit{Anfigonia:} fusione di due gameti aploidi per formare lo zigote diploide;
      \item \textit{Partenogenesi:} sviluppo dell'uovo non fecondato, senza variabilità genetica.
    \end{itemize}
\end{itemize}

\subsection{Metabolismo}

Il \textbf{metabolismo} è il complesso di trasformazioni chimiche che avvengono negli organismi viventi e si suddivide in:
\begin{itemize}
  \item \textbf{Anabolismo:} sintesi di molecole complesse a partire da molecole semplici, con consumo di energia (reazioni endoergoniche; es. sintesi proteica);
  \item \textbf{Catabolismo:} degradazione di molecole complesse in molecole semplici, con liberazione di energia (reazioni esoergoniche; es. respirazione cellulare, ciclo di Krebs).
\end{itemize}

Gli organismi si dividono in base alla fonte di carbonio:
\begin{itemize}
  \item \textbf{Autotrofi:} sintetizzano molecole organiche a partire da molecole inorganiche (fotosintesi);
  \item \textbf{Eterotrofi:} si nutrono di molecole organiche preformate da altri organismi.
\end{itemize}

\subsection{Rapporti interspecifici}

I rapporti tra organismi di specie diverse sono classificati come segue:

\begin{center}
\begin{tabular}{lll}
\toprule
\textbf{Tipo} & \textbf{Specie A} & \textbf{Specie B} \\
\midrule
Mutualismo & Beneficio ($+$) & Beneficio ($+$) \\
Commensalismo & Beneficio ($+$) & Neutrale (0) \\
Inquilinismo & Spazio condiviso & Neutrale (0) \\
Amensalismo & Neutrale (0) & Danno ($-$) \\
Parassitismo & Beneficio ($+$) & Danno ($-$) \\
\bottomrule
\end{tabular}
\end{center}

% ============================================================
\section{Citologia: strutture sub-cellulari}
% ============================================================

\subsection{Batteriofagi}

I \textbf{batteriofagi} (o fagi) sono virus che infettano i batteri, utilizzando la loro macchina metabolica per replicarsi. La replicazione può avvenire mediante:

\begin{enumerate}
  \item \textbf{Ciclo litico:} il fago inietta il proprio \dna{} nel batterio, che viene usato per produrre nuovi fagi; la cellula batterica si lisa rilasciando i nuovi virioni;
  \item \textbf{Ciclo lisogenico:} il \dna{} del fago si integra nel cromosoma batterico come \textit{profago} e si replica insieme ad esso senza lisare immediatamente la cellula.
\end{enumerate}

\subsection{Virus}

I \textbf{virus} sono entità biologiche acellulari, composti da:
\begin{itemize}
  \item \textbf{Genoma virale:} \dna{} o \rna{};
  \item \textbf{Capside:} involucro proteico formato da capsomeri, con funzioni protettive e di facilitazione dell'ingresso nella cellula ospite;
  \item \textbf{Pericapside} (o \textit{peplos}): involucro lipoproteico esterno al capside, non sempre presente.
\end{itemize}

I virus si classificano in base al tipo di cellula ospite (batteriofagi, virus vegetali, virus animali), al tipo di acido nucleico (\textit{desossiribovirus} a \dna, \textit{ribovirus} a \rna), e alla forma (elicoidale, allungata, sferica, poliedrica).

\subsection{Viroidi e Prioni}

I \textbf{viroidi} sono piccole molecole di \rna{} circolari, prive di capside, in grado di auto-replicarsi; infettano preferenzialmente le cellule vegetali.

I \textbf{prioni} sono proteine normalmente presenti nei tessuti sani. Quando la loro struttura tridimensionale si altera, le molecole misfoldate inducono il cambiamento conformazionale delle proteine native, innescando una reazione a catena responsabile di gravi malattie neurodegenerative.

\subsection{Cellula procariote}

Le cellule procariotiche si classificano in base a forma, temperatura ottimale di crescita, rapporto con l'ossigeno e concentrazione salina/acidità dell'ambiente.

\subsubsection{Struttura della cellula batterica}

\begin{itemize}
  \item \textbf{Capsula esterna}: presente in alcuni batteri, con pili e flagello;
  \item \textbf{Parete cellulare}: composta da membrana esterna e strato di peptidoglicano; la spessore dello strato di peptidoglicano distingue i batteri in:
    \begin{itemize}
      \item \textit{Gram-positivi:} strato spesso, si colorano di viola con la colorazione di Gram;
      \item \textit{Gram-negativi:} strato sottile, non si colorano;
    \end{itemize}
  \item \textbf{Membrana plasmatica};
  \item \textbf{Citoplasma}: soluzione acquosa contenente ribosomi (per la sintesi proteica);
  \item \textbf{DNA batterico} (nucleoide): libero nel citoplasma, circolare;
  \item \textbf{Plasmidi}: \dna{} circolare extracromosomico che conferisce caratteristiche aggiuntive (es. resistenza antibiotica).
\end{itemize}

% ============================================================
\section{Biotecnologie}
% ============================================================

\subsection{Principali tecniche di laboratorio}

\subsubsection{PCR (Polymerase Chain Reaction)}

La \textbf{PCR} è una tecnica che amplifica \textit{in vitro} specifiche sequenze di \dna. Il processo richiede: \dna{} stampo, due primer complementari ai filamenti della sequenza bersaglio, i quattro nucleotidi trifosfato e la \textit{Taq} polimerasi (un enzima termostabile). Il ciclo prevede:
\begin{enumerate}
  \item \textbf{Denaturazione} a 96\,°C: separazione delle due catene di \dna;
  \item \textbf{Annealing} a $\sim$50\,°C: appaiamento dei primer ai filamenti complementari;
  \item \textbf{Estensione} a 72\,°C: sintesi dei nuovi filamenti da parte della \textit{Taq} polimerasi.
\end{enumerate}
Ripetendo il ciclo (tipicamente 30--40 volte) si ottiene un'amplificazione esponenziale del frammento bersaglio.

\subsubsection{Elettroforesi su gel di agarosio}

Tecnica che separa frammenti di \dna{} (o proteine) in base alla dimensione, sfruttando un campo elettrico applicato a un gel di agarosio. I frammenti più corti migrano più velocemente verso l'elettrodo positivo. Il \dna{} è reso visibile per fluorescenza con un colorante intercalante (es. bromuro di etidio) sotto luce UV.

\subsubsection{Sequenziamento del DNA e Genoma umano}

Il sequenziamento permette di determinare l'ordine dei nucleotidi in una catena di \dna. Il genoma umano, completamente sequenziato nel 2003, è composto da $3 \times 10^9$ paia di basi, di cui soltanto il 3\% costituisce \dna{} codificante; il restante 97\% è rappresentato da sequenze non codificanti o intergeniche.

\subsubsection{Transgenesi e Clonazione}

La \textbf{transgenesi} consiste nell'introduzione di un gene esogeno (transgene) in un organismo ospite, con lo scopo di modificarne il fenotipo. La \textbf{clonazione} produce organismi con patrimonio genetico identico a quello di un individuo donatore, trasferendo il nucleo di una cellula somatica adulta in un ovocita enucleato.

Le \textbf{cellule staminali} sono cellule indifferenziate dotate di capacità di auto-rinnovamento e di differenziazione. Si classificano in:
\begin{itemize}
  \item \textit{Totipotenti:} possono generare qualsiasi tipo cellulare dell'organismo e i tessuti extraembrionali;
  \item \textit{Pluripotenti/multipotenti:} possono generare molti tipi cellulari, ma non tutti.
\end{itemize}

% ============================================================
\chapter{Biologia Cellulare}
% ============================================================

\begin{tcolorbox}[summarybox, title={Sommario del capitolo}]
Questo capitolo descrive l'organizzazione strutturale e funzionale della cellula eucariote: la membrana plasmatica (struttura, proteine, meccanismi di trasporto), il citoscheletro, le giunzioni cellulari, il sistema delle endomembrane (reticolo endoplasmatico, complesso di Golgi, lisosomi, perossisomi, vacuoli) e gli organuli per la produzione di energia (mitocondri e cloroplasti).
\end{tcolorbox}

% ============================================================
\section{Membrana plasmatica}
% ============================================================

La \textbf{membrana plasmatica} delimita la cellula e regola il passaggio di sostanze tra l'interno e l'esterno. Ha una struttura \textbf{trilaminare} organizzata come un \textit{doppio strato fosfolipidico}, con le teste idrofile rivolte verso i comparti acquosi (interno ed esterno alla cellula) e le code idrofobiche degli acidi grassi orientate verso l'interno della membrana. Nel doppio strato sono intercalati:
\begin{itemize}
  \item \textbf{Colesterolo:} anfipatico come i fosfolipidi, aumenta la rigidità e la fluidità della membrana a seconda della temperatura;
  \item \textbf{Glicolipidi:} lipidi con catene saccaridiche rivolte verso l'esterno cellulare.
\end{itemize}

\subsection{Modelli di struttura della membrana}

Il \textbf{modello di Davson--Danielli} (1954) prevedeva uno strato proteico monomolecolare su entrambe le superfici del doppio strato.

Il \textbf{modello a mosaico fluido di Singer e Nicolson} (1972), tuttora valido, descrive la membrana come un sistema dinamico in cui le proteine sono inserite o associate al doppio strato lipidico, libero di scorrere lateralmente:
\begin{itemize}
  \item \textit{Proteine periferiche} (estrinseche): associate alla superficie interna o esterna del doppio strato;
  \item \textit{Proteine integrali} (intrinseche): ancorate alla porzione idrofobica centrale;
  \item \textit{Proteine transmembrana:} attraversano interamente la membrana; possono essere \textit{single pass} (una sola traversata) o \textit{multiple pass} (più attraversamenti).
\end{itemize}

La fluidità del doppio strato è stata dimostrata dall'\textbf{esperimento di Frye--Edidin}: cellule umane (con proteine di membrana marcate in rosso) e cellule di topo (marcate in verde) vengono fuse in ibridomi; dopo 40 minuti di incubazione, le proteine si distribuiscono uniformemente su tutta la superficie, a dimostrazione della mobilità laterale dei fosfolipidi (e delle proteine ad essi associate).

\subsection{Funzioni della membrana plasmatica}

\begin{enumerate}
  \item Protezione meccanica e chimica della cellula;
  \item Ancoraggio alla matrice extracellulare;
  \item Mantenimento della forma cellulare;
  \item Regolazione degli scambi di sostanze nutritive;
  \item Trasporto attivo e passivo di sostanze;
  \item Risposta agli stimoli ormonali tramite recettori specifici;
  \item Riconoscimento cellulare e comunicazione intercellulare;
  \item Proprietà antigeniche e risposta immunitaria.
\end{enumerate}

\subsection{Meccanismi di trasporto}

I meccanismi di trasporto attraverso la membrana si distinguono in \textbf{passivi} (secondo gradiente, senza consumo di energia) e \textbf{attivi} (contro gradiente, con consumo di \atp).

\begin{center}
\begin{tabular}{lll}
\toprule
\textbf{Meccanismo} & \textbf{Direzione} & \textbf{Energia} \\
\midrule
Diffusione semplice & Secondo gradiente & No \\
Diffusione facilitata (canale) & Secondo gradiente & No \\
Diffusione facilitata (carrier) & Secondo gradiente & No \\
Trasporto attivo & Contro gradiente & Sì (1 ATP) \\
Esocitosi & Verso l'esterno & Sì \\
Fagocitosi & Verso l'interno & Sì \\
Pinocitosi & Verso l'interno & Sì \\
Endocitosi mediata da recettore & Verso l'interno & Sì \\
\bottomrule
\end{tabular}
\end{center}

\begin{tcolorbox}[exbox, title={Pompa sodio-potassio}]
  La pompa Na$^+$/K$^+$-ATPasi è un esempio paradigmatico di trasporto attivo primario: per ogni molecola di ATP idrolizzata, espelle 3 ioni Na$^+$ verso l'esterno e introduce 2 ioni K$^+$ all'interno della cellula, mantenendo i gradienti ionici essenziali per l'eccitabilità cellulare.
\end{tcolorbox}

% ============================================================
\section{Il citoscheletro}
% ============================================================

Il \textbf{citoscheletro} è una rete dinamica di filamenti proteici che pervade il citoplasma, ancorata alla membrana plasmatica tramite proteine specializzate. Svolge funzioni strutturali, di trasporto e di motilità cellulare. È composto da tre tipi di filamenti:

\subsection{Microtubuli}

I \textbf{microtubuli} sono filamenti cilindrici cavi con diametro di 20--25\,nm, formati da dimeri di $\alpha$-tubulina e $\beta$-tubulina organizzati in anelli di 13 protofilamenti. Presentano polarità intrinseca: i dimeri vengono aggiunti all'estremità positiva ($+$) e rimossi da quella negativa ($-$).

Funzioni dei microtubuli:
\begin{itemize}
  \item Impalcatura per la forma cellulare;
  \item Guida per il movimento di organuli e vescicole (tramite motori molecolari: \textit{chinesina} verso l'estremità $+$; \textit{dineina} verso l'estremità $-$);
  \item Formazione del fuso mitotico, essenziale per la separazione dei cromosomi;
  \item Costituzione dell'assonema di ciglia e flagelli (struttura 9+2).
\end{itemize}

I \textbf{centrioli} sono organuli cilindrici formati da 9 triplette di microtubuli; si trovano sempre in coppia perpendicolare (il \textit{centrosoma}) e durante la divisione cellulare migrano ai poli opposti, nucleando i microtubuli del fuso mitotico.

\subsection{Filamenti intermedi}

I \textbf{filamenti intermedi} hanno un diametro di 8--12\,nm e conferiscono resistenza meccanica alla cellula. Sono formati da subunità proteiche con struttura ad $\alpha$-elica centrale e domini terminali globulari, che si associano in dimeri, tetrameri e unità filamentose più lunghe.

Nella cellula i filamenti intermedi sono ancorati sia al nucleo sia alla membrana plasmatica tramite i desmosomi, garantendo la coesione meccanica tra cellule adiacenti.

\subsection{Microfilamenti}

I \textbf{microfilamenti} (diametro 5--7\,nm) sono polimeri flessibili di actina-G globulare, che polimerizzano (con consumo di ATP) in filamenti ad elica di actina-F; due filamenti si associano a formare i microfilamenti definitivi.

Diverse proteine regolatrici controllano la dinamica dei microfilamenti: proteine di nucleazione, sequestro, incappucciamento, depolimerizzazione e di legame alla membrana. I microfilamenti compongono i \textbf{microvilli} dell'epitelio intestinale, aumentando notevolmente la superficie di assorbimento.

% ============================================================
\section{Le giunzioni cellulari}
% ============================================================

Le \textbf{giunzioni cellulari} sono strutture specializzate che connettono le cellule tra loro e alla matrice extracellulare.

\begin{center}
\begin{tabular}{lll}
\toprule
\textbf{Tipo di giunzione} & \textbf{Funzione principale} & \textbf{Comunicazione} \\
\midrule
GAP (comunicanti) & Scambio di ioni e piccole molecole & Sì \\
Aderenti & Coesione meccanica & No \\
Strette & Barriera impermeabile & No \\
Desmosomi & Ancoraggio meccanico & No \\
Plasmodesmi (vegetali) & Passaggio di acqua e ioni & Sì \\
\bottomrule
\end{tabular}
\end{center}

\subsection{Giunzioni GAP}

Sono canali proteici (formati da \textit{connessine}) che attraversano le membrane di due cellule adiacenti, creando canali idrofili regolabili. Permettono la comunicazione metabolica ed elettrica rapida tra cellule (fondamentale nel miocardio).

\subsection{Desmosomi e Giunzioni Strette}

I \textbf{desmosomi} ancorizzano la cheratina del citoscheletro di cellule adiacenti, formando una struttura densa (la \textit{placca}) altamente resistente alle sollecitazioni meccaniche. Sono caratteristici del tessuto epiteliale.

Le \textbf{giunzioni strette} formano cinture di fusione tra le membrane di cellule adiacenti, abolendo lo spazio intercellulare e impedendo la diffusione paracellulare di sostanze (es. barriera intestinale, ematoencefalica).

% ============================================================
\section{Le endomembrane}
% ============================================================

\subsection{Reticolo endoplasmatico}

Il \textbf{reticolo endoplasmatico} (RE) è una rete di membrane che si estende nel citoplasma, in continuità con l'involucro nucleare. Si divide in:
\begin{itemize}
  \item \textbf{RE liscio:} sistema di tubuli membranosi privo di ribosomi; svolge funzioni di sintesi di lipidi e ormoni steroidei, detossificazione epatica e sequestro di ioni Ca$^{2+}$;
  \item \textbf{RE rugoso:} cisterne appiattite ricoperte da ribosomi; è il sito principale di sintesi, ripiegamento e glicosilazione delle proteine destinate alla secrezione, alla membrana o a organuli.
\end{itemize}

\subsection{Complesso di Golgi}

Il \textbf{complesso di Golgi} è interposto tra il RE e la membrana plasmatica. È formato da pile di cisterne appiattite (la \textit{dictiocinetica}), organizzate in tre regioni funzionali:
\begin{itemize}
  \item Cisterne \textbf{cis} (rivolte al RE): ricevono le proteine dal RE;
  \item Cisterne \textbf{mediali}: sede di modificazioni post-traduzionali;
  \item Cisterne \textbf{trans} (rivolte alla membrana): smistano le proteine alla loro destinazione finale.
\end{itemize}

Le modificazioni effettuate dal Golgi includono: rimozione di porzioni della catena amminoacidica, aggiunta di gruppi funzionali, glicosilazione e lipidazione.

\subsection{Lisosomi}

I \textbf{lisosomi} sono organuli digestivi delimitati da membrana, contenenti circa 50 idrolasi acide (proteasi, lipasi, nucleasi, ecc.) attive a pH 4--5, mantenuto da una pompa protonica. Svolgono:
\begin{enumerate}
  \item Degradazione di materiale esterno ingerito per fagocitosi;
  \item \textit{Turnover} degli organuli cellulari (\textit{autofagia}): l'organulo obsoleto viene circondato da membrana del RE, formando un autofagosoma che si fonde con il lisosoma.
\end{enumerate}

\subsection{Perossisomi e Vacuoli}

I \textbf{perossisomi} (diametro 0,1--1,0\,$\mu$m) contengono ossidasi e catalasi: le prime generano $\text{H}_2\text{O}_2$ ossidando substrati organici; la catalasi decompone rapidamente $\text{H}_2\text{O}_2$ in acqua e ossigeno, neutralizzando così la sua tossicità.

I \textbf{vacuoli} occupano circa il 90\% del volume delle cellule vegetali mature. Sono delimitati dal \textit{tonoplasto} (membrana con pompa protonica che mantiene pH acido). Il turgore del vacuolo, determinato dall'ingresso di acqua per osmosi, conferisce rigidità alla cellula vegetale.

% ============================================================
\section{Mitocondri e Cloroplasti}
% ============================================================

\subsection{Mitocondri}

I \textbf{mitocondri} sono gli organuli respiratori delle cellule eucariote, responsabili della produzione della maggior parte dell'\atp. Sono formati da:
\begin{enumerate}
  \item \textbf{Membrana esterna:} liscia e permeabile a piccole molecole;
  \item \textbf{Spazio intermembrana};
  \item \textbf{Membrana interna:} ripiegata a formare le \textit{creste}, sede degli enzimi della catena respiratoria e dell'ATP sintasi;
  \item \textbf{Matrice mitocondriale:} sede degli enzimi del ciclo di Krebs; contiene anche \dna{} mitocondriale, ribosomi, tRNA e mRNA.
\end{enumerate}

\subsubsection{Respirazione cellulare}

La produzione di \atp{} a partire dal glucosio avviene in tre stadi:

\begin{enumerate}
  \item \textbf{Glicolisi} (citoplasma): il glucosio è convertito in 2 molecole di piruvato, con produzione netta di 2 \atp{} e 2 \nadh;

  \item \textbf{Ciclo di Krebs} (matrice mitocondriale): l'acetil-CoA (derivato dall'ossidazione del piruvato) viene completamente ossidato a \CO, con produzione di \nadh, FADH$_2$ e 2 \atp{} per ciclo;

  \item \textbf{Fosforilazione ossidativa} (membrane delle creste):
    \begin{itemize}
      \item \textit{Catena di trasporto degli elettroni:} gli elettroni del \nadh{} e del FADH$_2$ vengono trasferiti lungo una serie di complessi proteici trans-membrana, generando un gradiente protonico ($\Delta\mu_{\text{H}^+}$) tra lo spazio intermembrana e la matrice;
      \item \textit{ATP sintasi:} sfrutta il gradiente protonico per fosforilare l'ADP, producendo $\sim$32 \atp.
    \end{itemize}
\end{enumerate}

\begin{tcolorbox}[notebox, title={Origine endosimbiontica dei mitocondri}]
  L'ipotesi endosimbiontica suggerisce che i mitocondri derivino da batteri aerobi anticamente inglobati da cellule eucariotiche antenate. A sostegno di questa teoria: il \dna{} mitocondriale è circolare (come quello batterico), i ribosomi mitocondriali sono di tipo 70S (batterico) e la membrana interna (le creste) corrisponde alla membrana originale del batterio endosimbionte.
\end{tcolorbox}

\subsection{Cloroplasti}

I \textbf{cloroplasti} sono organuli esclusivi delle cellule vegetali (e di alcune alghe e batteri fotosintetici). La loro funzione è la \textbf{fotosintesi clorofilliana}: la conversione di energia luminosa in energia chimica immagazzinata nel glucosio.

La struttura comprende:
\begin{itemize}
  \item Membrana esterna liscia (con porine);
  \item Membrana interna (con trasportatori), ripiegata a formare i \textit{tilacoidi} (impilati in \textit{grana});
  \item \textbf{Stroma:} spazio interno contenente gli enzimi per la sintesi dei carboidrati, \dna{} residuo (circolare) e ribosomi.
\end{itemize}

\subsubsection{Fotosintesi clorofilliana}

La fotosintesi avviene in due fasi:

\begin{enumerate}
  \item \textbf{Fase luminosa} (membrana dei tilacoidi): l'energia luminosa è assorbita dai pigmenti fotosintetici (clorofille) e convertita in energia chimica (\atp{} e \nadph), con produzione di \Oc{} per fotolisi dell'acqua:
  \begin{equation}
    6\,\HO + \text{luce} \;\xrightarrow{\text{Fotosistemi I \& II}}\; 6\,\Oc + \text{ATP} + \text{NADPH}
  \end{equation}

  \item \textbf{Fase oscura} (ciclo di Calvin, nello stroma): l'\atp{} e il \nadph{} prodotti nella fase luminosa vengono utilizzati per ridurre la \CO{} a glucosio:
  \begin{equation}
    6\,\CO + 12\,\HO + \text{ATP} + \text{NADPH} \;\to\; \text{C}_6\text{H}_{12}\text{O}_6 + 6\,\HO + 6\,\Oc
  \end{equation}
\end{enumerate}

% ============================================================
\chapter{Ciclo Cellulare e Genetica}
% ============================================================

\begin{tcolorbox}[summarybox, title={Sommario del capitolo}]
Questo capitolo descrive la struttura del nucleo e la duplicazione del \dna, il ciclo cellulare (interfase, mitosi, meiosi), la sintesi proteica (trascrizione e traduzione), il codice genetico, le leggi di Mendel e le alterazioni cromosomiche.
\end{tcolorbox}

% ============================================================
\section{Nucleo e DNA}
% ============================================================

\subsection{Struttura del nucleo}

Il \textbf{nucleo} è il centro di controllo della cellula eucariote, sede del materiale genetico. È delimitato dall'\textbf{involucro nucleare}, formato da due membrane distanti 10--50\,nm:
\begin{itemize}
  \item \textbf{Membrana esterna}: in continuità con il RE rugoso;
  \item \textbf{Membrana interna}: ancorata alla \textit{lamina nucleare}, una rete fibrillare di proteine di tipo filamento intermedio.
\end{itemize}

I \textbf{pori nucleari} attraversano entrambe le membrane; hanno struttura ottagonale (8 ripetizioni di 30 tipi proteici diversi) e regolano il traffico bidirezionale di macromolecole tra nucleo e citoplasma. Le proteine con segnale di localizzazione nucleare (NLS) vengono riconosciute dalle importine e trasportate attivamente nel nucleo in un processo che richiede la proteina GTPasica Ran.

\subsection{Cromosomi}

I \textbf{cromosomi} sono la forma condensata della cromatina, visibile durante la divisione cellulare. La cromatina è formata da \dna{} (carico negativamente) associato a istoni (cariche positive); può esistere in due stati:
\begin{itemize}
  \item \textbf{Eucromatina:} forma dispersa, trascrizionalmente attiva;
  \item \textbf{Eterocromatina:} forma condensata, trascrizionalmente silenziosa.
\end{itemize}

L'organizzazione del \dna{} in cromosomi segue livelli successivi di compattamento:
\begin{enumerate}
  \item 146 paia di basi di \dna{} si avvolgono attorno a 8 istoni formando il \textit{nucleosoma} ($\phi = 10$\,nm);
  \item I nucleosomi si compattano tramite l'istone H1 formando una fibra da 30\,nm;
  \item Ulteriori superavvolgimenti portano alla formazione del cromosoma metafasico.
\end{enumerate}

Ogni cromosoma è composto da:
\begin{itemize}
  \item Un'unica molecola di \dna{} continua;
  \item \textbf{Telomeri:} sequenze ripetute alle estremità;
  \item \textbf{Due cromatidi fratelli:} unità parallele identiche;
  \item \textbf{Centromero:} regione di constrizione primaria, sito di attacco per i microtubuli del fuso (tramite i cinetocori).
\end{itemize}

\subsection{Duplicazione del DNA}

La duplicazione del \dna{} è un processo semiconservativo che coinvolge diversi enzimi specializzati:

\begin{enumerate}
  \item La \textbf{DNA girasi} (topoisomerasi II) elimina la tensione torsionale a monte della forca replicativa;
  \item La \textbf{DNA elicasi} svolge e separa i due filamenti;
  \item Le \textbf{proteine SSB} stabilizzano i filamenti singoli;
  \item La \textbf{RNA primasi} sintetizza i primer (innesco) di RNA;
  \item La \textbf{DNA polimerasi} estende i primer, sintetizzando il nuovo filamento in direzione $5'\!\to\!3'$. Sul filamento continuo (\textit{leading strand}) la sintesi è continua; sul filamento discontinuo (\textit{lagging strand}) vengono prodotti i frammenti di Okazaki, successivamente uniti dalla \textbf{DNA ligasi}.
\end{enumerate}

\begin{tcolorbox}[notebox, title={Errori di replicazione}]
  La DNA polimerasi possiede un'attività esonucleasica correttrice di bozze che riduce gli errori di incorporazione. Qualora l'elica usata come stampo sia essa stessa mutata, l'errore si trasmette alla progenie come \textbf{mutazione}.
\end{tcolorbox}

% ============================================================
\section{Ciclo cellulare}
% ============================================================

Il \textbf{ciclo cellulare} è la sequenza ordinata di eventi che porta alla duplicazione e alla divisione di una cellula. Si articola in:

\begin{itemize}
  \item \textbf{Interfase} ($\sim$23 ore in una cellula mammiferia tipica):
    \begin{itemize}
      \item \textbf{Fase G$_0$:} quiescenza (arresto reversibile del ciclo);
      \item \textbf{Fase G$_1$:} crescita cellulare e svolgimento delle funzioni normali;
      \item \textbf{Fase S:} duplicazione del \dna{} e dei cromosomi;
      \item \textbf{Fase G$_2$:} preparazione alla mitosi, duplicazione dei centrioli;
    \end{itemize}
  \item \textbf{Fase M:}
    \begin{itemize}
      \item \textbf{Mitosi:} segregazione equa dei cromosomi nelle due cellule figlie;
      \item \textbf{Citocinesi:} divisione fisica del citoplasma.
    \end{itemize}
\end{itemize}

\subsection{Mitosi}

La \textbf{mitosi} assicura la distribuzione equa del materiale genetico duplicato alle due cellule figlie e si divide in cinque fasi:

\begin{enumerate}
  \item \textbf{Profase:} condensazione della cromatina in cromosomi; smontaggio del citoscheletro e assemblaggio del fuso mitotico; frammentazione del Golgi e del RE; dissoluzione dell'involucro nucleare;
  \item \textbf{Prometafase:} i microtubuli cromosomici si agganciano ai cinetocori; i cromosomi si spostano verso l'equatore cellulare;
  \item \textbf{Metafase:} allineamento dei cromosomi sulla \textit{piastra metafasica}, connessi ai poli tramite microtubuli cromosomici;
  \item \textbf{Anafase:} separazione dei cromatidi fratelli per clivaggio del centromero; i cromosomi migrano verso i poli opposti; i poli del fuso si allontanano;
  \item \textbf{Telofase:} i cromosomi decondensano; si riformano l'involucro nucleare, il Golgi e il RE; inizia la citocinesi (formazione di un solco equatoriale con contrazione dei microfilamenti di actina).
\end{enumerate}

\subsection{Meiosi}

La \textbf{meiosi} è il processo di divisione cellulare che genera i gameti, cellule aploidi (con metà del corredo cromosomico) attraverso due divisioni successive (meiosi I e meiosi II) precedute da una sola duplicazione del \dna.

\begin{tcolorbox}[defbox, title={Crossing over}]
  Durante la proafase I, i cromosomi omologhi si appaiano formando le \textit{tetradi} e possono scambiare segmenti di \dna{} per ricombinazione omologa (\textbf{crossing over}), aumentando la variabilità genetica dei gameti.
\end{tcolorbox}

Le tappe della meiosi sono:
\begin{enumerate}
  \item \textbf{Profase I}: condensazione dei cromosomi, formazione delle tetradi, possibilità di crossing over;
  \item \textbf{Metafase I}: le tetradi si allineano sul piano equatoriale;
  \item \textbf{Anafase I}: separazione dei cromosomi omologhi verso i poli opposti (i cromatidi fratelli restano uniti);
  \item \textbf{Telofase I e citocinesi}: formazione di due cellule aploidi;
  \item \textbf{Profase II -- Metafase II -- Anafase II -- Telofase II}: le due cellule figlie si dividono ulteriormente, separando i cromatidi, per un totale di 4 cellule aploidi (gameti/spore).
\end{enumerate}

% ============================================================
\section{Sintesi proteica}
% ============================================================

La \textbf{sintesi proteica} avviene in due fasi successive:
\begin{enumerate}
  \item \textbf{Trascrizione} (nel nucleo): una sequenza di \dna{} è copiata in un filamento di mRNA;
  \item \textbf{Traduzione} (nel citoplasma, sui ribosomi): l'mRNA è decodificato dai tRNA e dai ribosomi per produrre una catena polipeptidica.
\end{enumerate}

\subsection{Trascrizione}

La \textbf{trascrizione} è catalizzata dalla RNA polimerasi, che prende come stampo il filamento di \dna{} $3'\!\to\!5'$, sintetizzando un filamento di mRNA nella direzione $5'\!\to\!3'$.

Le fasi principali sono:
\begin{enumerate}
  \item \textbf{Inizio:} la RNA polimerasi riconosce la sequenza promotore e inizia la sintesi;
  \item \textbf{Allungamento:} la RNA polimerasi svvolge il \dna{} localmente, transcrive e riavvolge;
  \item \textbf{Terminazione:} la polimerasi trascrive sequenze ricche di C-G che formano strutture a forcina nell'mRNA nascente, bloccando il progresso della polimerasi stessa.
\end{enumerate}

Il trascritto primario (pre-mRNA) subisce modifiche post-trascrizionali:
\begin{itemize}
  \item Aggiunta del \textbf{cappuccio 5' (CAP)}: metil-guanosina che protegge l'estremità 5' dalla degradazione e facilita il riconoscimento ribosomale;
  \item Aggiunta della \textbf{coda poli(A)}: 200--300 residui di adenina all'estremità 3', che proteggono dall'esonucleasi;
  \item \textbf{Splicing}: rimozione delle sequenze non codificanti (introni) e giunzione degli esoni.
\end{itemize}

\subsection{Traduzione}

La \textbf{traduzione} avviene sui ribosomi (complessi 80S negli eucarioti, formati da una subunità maggiore 60S e una minore 40S) e si articola in tre fasi:

\begin{enumerate}
  \item \textbf{Inizio:} l'mRNA si associa alla subunità minore del ribosoma; il tRNA iniziatore (portante metionina) riconosce il codone di inizio AUG; si assembla il ribosoma completo;
  \item \textbf{Allungamento:} ciclo ripetitivo in cui un aminoacil-tRNA entra nel sito A, si forma il legame peptidico con la catena in sito P, il ribosoma trasla di un codone (translocazione);
  \item \textbf{Terminazione:} un codone di stop (UAA, UAG o UGA) viene riconosciuto da un fattore di rilascio; la catena polipeptidica viene rilasciata e il ribosoma si dissocia.
\end{enumerate}

\subsection{Codice genetico}

Il \textbf{codice genetico} è l'insieme delle 64 triplette di basi (codoni) che specificano i 20 amminoacidi proteici o i segnali di stop. Possiede quattro proprietà fondamentali:

\begin{enumerate}
  \item \textbf{Degenerato} (o ridondante): più codoni possono codificare lo stesso amminoacido (es. leucina è codificata da 6 codoni);
  \item \textbf{Univoco:} ogni codone specifica un solo amminoacido;
  \item \textbf{Non sovrapposto:} ogni base appartiene a un solo codone;
  \item \textbf{Universale:} quasi tutti gli organismi viventi usano lo stesso codice (con rare eccezioni nei mitocondri e in alcuni protozoi).
\end{enumerate}

% ============================================================
\section{Genetica di Mendel}
% ============================================================

Gregor Mendel (1822--1884) condusse esperimenti sistematici su \textit{Pisum sativum} (pisello), sfruttandone la facilità di coltura, il breve ciclo vitale e la presenza di entrambi gli apparati riproduttivi sullo stesso individuo. Formulò tre leggi fondamentali:

\subsection{Prima legge: Legge della dominanza}

Incrociando due linee pure per un carattere alternativo, tutti gli individui della generazione F$_1$ manifestano un solo fenotipo, detto \textbf{dominante}; il fenotipo \textbf{recessivo} non si esprime in F$_1$.

\subsection{Seconda legge: Legge della segregazione}

Nella formazione dei gameti, i due alleli di un gene \textbf{segregano} (si separano) in modo che ogni gamete riceva un solo allele. Pertanto, in F$_2$, il rapporto fenotipico tra individui dominanti e recessivi è \textbf{3:1}.

\subsection{Terza legge: Legge dell'indipendenza}

Gli alleli che determinano caratteri \textbf{diversi} segregano in modo indipendente l'uno dall'altro, producendo in F$_2$ tutte le combinazioni possibili in rapporto 9:3:3:1. Questa legge vale per geni su cromosomi diversi (non in linkage).

\subsection{Eccezioni alle leggi di Mendel}

\begin{itemize}
  \item \textbf{Dominanza incompleta:} gli eterozigoti mostrano un fenotipo intermedio tra i due omozigoti (es. colore dei fiori della digitale);
  \item \textbf{Codominanza:} entrambi gli alleli si esprimono ugualmente (es. gruppo sanguigno AB del sistema ABO);
  \item \textbf{Epistasi:} l'espressione di un gene è soppressa dall'azione di un altro gene (es. colore del pelo nel labrador retriever).
\end{itemize}

% ============================================================
\section{Cariotipo e alterazioni cromosomiche}
% ============================================================

Il \textbf{cariotipo} è l'assetto cromosomico caratteristico di una specie, visualizzabile durante la metafase della mitosi. Il cariotipo umano consta di 46 cromosomi (23 coppie): 22 coppie di \textbf{autosomi} e 1 coppia di \textbf{cromosomi sessuali} (XX nelle femmine, XY nei maschi).

\subsection{Principali alterazioni cromosomiche}

\begin{itemize}
  \item \textbf{Delezione:} perdita di un segmento cromosomico;
  \item \textbf{Duplicazione:} inserimento di un segmento duplicato nel cromosoma omologo;
  \item \textbf{Traslocazione:} scambio reciproco di segmenti tra cromosomi non omologhi (es. Linfoma di Burkitt);
  \item \textbf{Inversione:} riinserzione di un segmento ruotato di 180°.
\end{itemize}

\subsection{Non disgiunzione meiotica}

Quando una coppia di omologhi non si separa correttamente durante la meiosi, si produce \textbf{aneuploidia}. La fecondazione di un gamete normale ($n$) con uno abnorme ($n+1$ o $n-1$) origina:
\begin{itemize}
  \item \textbf{Trisomia 21} (Sindrome di Down): tre cromosomi 21; $2n+1 = 47$ cromosomi;
  \item \textbf{Sindrome di Turner} (45, X0): fenotipo femminile, ovaie iposviluppate, sterilità;
  \item \textbf{Sindrome di Klinefelter} (47, XXY): fenotipo maschile, testicoli iposviluppati, sterilità.
\end{itemize}

\subsection{Mutazioni puntiformi}

Le \textbf{mutazioni puntiformi} interessano una singola coppia di basi e si classificano in:
\begin{enumerate}
  \item \textbf{Mutazione silente:} cambio di base che non modifica l'amminoacido (per degenerazione del codice);
  \item \textbf{Mutazione missenso:} cambio di base con sostituzione di un amminoacido (es. anemia falciforme: la sostituzione di acido glutammico con valina in posizione 6 dell'emoglobina);
  \item \textbf{Mutazione non senso:} creazione di un codone di stop prematuro, con troncatura della proteina;
  \item \textbf{Mutazione frameshift:} inserzione o delezione di 1 o 2 basi che altera il quadro di lettura.
\end{enumerate}

% ============================================================
\chapter{Istologia}
% ============================================================

\begin{tcolorbox}[summarybox, title={Sommario del capitolo}]
L'istologia studia i \textbf{tessuti}: insiemi di cellule specializzate che cooperano per svolgere funzioni definite. I quattro tessuti fondamentali dell'organismo animale sono: epiteliale, connettivo, muscolare e nervoso.
\end{tcolorbox}

% ============================================================
\section{Tessuto epiteliale}
% ============================================================

Il \textbf{tessuto epiteliale} riveste le superfici esterne e interne del corpo, delimita le vie di transito dei fluidi corporei e costituisce le ghiandole. Le sue caratteristiche principali sono:
\begin{itemize}
  \item \textbf{Cellularità elevata:} cellule molto vicine, con scarsa matrice extracellulare;
  \item \textbf{Polarità:} le cellule presentano un polo apicale (verso il lume o la superficie) e un polo basale (verso la lamina basale);
  \item \textbf{Posizione sulla lamina basale:} tutte le cellule epiteliali poggiano su una lamina basale, tramite emidesmosomi;
  \item \textbf{Non vascolarizzato:} le cellule si nutrono per diffusione dal tessuto connettivo sottostante;
  \item \textbf{Elevata capacità proliferativa:} necessaria per il ricambio continuo delle cellule.
\end{itemize}

\subsection{Classificazione morfologica}

\begin{center}
\begin{tabular}{lll}
\toprule
\textbf{Forma cellulare} & \textbf{Strato singolo} & \textbf{Pluristratificato} \\
\midrule
Piatta/pavimentosa & Semplice & Stratificato \\
Cubica & Semplice & Stratificato \\
Prismatica/cilindrica & Semplice & Stratificato \\
\bottomrule
\end{tabular}
\end{center}

Oltre a queste categorie, esistono:
\begin{itemize}
  \item \textbf{Epitelio pseudostratificato:} un solo strato cellulare con nuclei a livelli differenti; tutte le cellule poggiano sulla lamina basale, ma non tutte raggiungono la superficie apicale;
  \item \textbf{Epitelio di transizione (urotelio):} cambia forma a seconda dello stato di distensione dell'organo (es. vescica urinaria).
\end{itemize}

\subsection{Epidermide}

L'epidermide è l'epitelio pavimentoso stratificato cheratinizzato che riveste la cute. È composta da cheratinociti, cellule che vanno incontro ad \textit{apoptosi} differenziativa (morte programmata con accumulo di cheratina). La struttura è organizzata in strati (dal più profondo al più superficiale):

\begin{enumerate}
  \item \textbf{Strato basale:} cellule cubiche/cilindriche con intensa attività proliferativa; unico strato in contatto con la lamina basale;
  \item \textbf{Strato spinoso:} cellule poliedriche con prolungamenti (spine) che formano desmosomi;
  \item \textbf{Strato granuloso:} cellule appiattite con granuli di cheratoialina; prime evidenze di apoptosi;
  \item \textbf{Strato lucido:} cellule appiattite con filamenti di cheratina paralleli (presente solo nella cute spessa);
  \item \textbf{Strato corneo:} lamelle cornee anucleate, costituite quasi esclusivamente da cheratina; desquamano continuamente.
\end{enumerate}

Altre cellule dell'epidermide:
\begin{itemize}
  \item \textbf{Melanociti} (strato basale/spinoso): producono melanina, responsabile della pigmentazione cutanea e della protezione dalle radiazioni UV;
  \item \textbf{Cellule di Langerhans} (strato spinoso): cellule dendritiche presentanti l'antigene, coinvolte nella risposta immunitaria cutanea;
  \item \textbf{Cellule di Merkel} (strato basale): meccanocettori.
\end{itemize}

\subsection{Epitelio ghiandolare}

Le \textbf{ghiandole} sono strutture specializzate nella produzione e secrezione di sostanze specifiche (\textit{secreti}):
\begin{itemize}
  \item \textbf{Ghiandole esocrine:} riversano i secreti in cavità o sulla superficie corporea tramite dotti escretori (es. ghiandole salivari, sebacee);
  \item \textbf{Ghiandole endocrine:} riverano i secreti (ormoni) direttamente nel sangue o nel liquido interstiziale (es. tiroide, surrene).
\end{itemize}

Le ghiandole esocrine si classificano in base alla modalità di secrezione:
\begin{itemize}
  \item \textbf{Merocrine/eccrine:} secrezione per esocitosi, senza perdita di citoplasma;
  \item \textbf{Apocrine:} secrezione con perdita della parte apicale del citoplasma;
  \item \textbf{Olocrine:} secrezione per distruzione completa della cellula (es. ghiandole sebacee).
\end{itemize}

% ============================================================
\section{Tessuto connettivo}
% ============================================================

Il \textbf{tessuto connettivo} comprende un ampio insieme di tessuti distribuiti in tutto il corpo, che non entrano mai in contatto diretto con l'esterno e svolgono funzioni di sostegno strutturale, riempimento degli spazi, trasporto e difesa. A differenza del tessuto epiteliale, è:
\begin{itemize}
  \item Ricco di matrice extracellulare;
  \item Vascolarizzato;
  \item Privo di polarità;
  \item Non poggiato su lamina basale.
\end{itemize}

\subsection{Classificazione}

\begin{enumerate}
  \item \textbf{Tessuto connettivo propriamente detto:} fibroso lasso, fibroso denso, elastico, reticolare, adiposo;
  \item \textbf{Tessuto di sostegno:} cartilagineo, osseo;
  \item \textbf{Tessuto emolinfopoietico:} sangue e organi ematopoietici.
\end{enumerate}

\subsection{Componenti cellulari del tessuto connettivo}

\begin{itemize}
  \item \textbf{Fibroblasti:} cellule affusolate che producono le componenti della matrice extracellulare (fibre e sostanza fondamentale); a riposo diventano fibrociti;
  \item \textbf{Mastociti:} liberano eparina (anticoagulante) e istamina (vasodilatatore); coinvolti nelle reazioni allergiche;
  \item \textbf{Macrofagi:} cellule fagocitiche che eliminano patogeni, detriti cellulari e sostanze estranee;
  \item \textbf{Adipociti:} cellule con grande gocciola lipidica intracitoplasmatica; costituiscono il tessuto adiposo;
  \item \textbf{Linfociti, plasmacellule, granulociti:} leucociti migrati dal sangue con funzioni di difesa immunitaria.
\end{itemize}

\subsection{Matrice extracellulare}

La matrice extracellulare è composta da:
\begin{enumerate}
  \item \textbf{Sostanza fondamentale:} gel idratato contenente proteoglicani (catene di glicosamminoglicani, GAG, ancorate a un asse proteico) e glicoproteine adesive (fibronectina, laminina, osteonectina);
  \item \textbf{Componente fibrillare:}
    \begin{itemize}
      \item \textit{Fibre collagene:} molto resistenti alla trazione; formate da tropocollagene ($\to$ microfibrille $\to$ fibrille $\to$ fibra collagene);
      \item \textit{Fibre reticolari:} fibre sottili ramificate a formare reticoli tridimensionali flessibili;
      \item \textit{Fibre elastiche:} nucleo di elastina con guaina di microfibrille di fibrillina; permettono al tessuto di allungarsi fino al 150\% della lunghezza di riposo.
    \end{itemize}
\end{enumerate}

\subsection{Tessuto connettivo propriamente detto}

\subsubsection{Tessuto connettivo fibroso lasso}
Costituisce il ``tessuto d'imballaggio'' del corpo; abbondante sostanza fondamentale, poche fibre, scarsi fibroblasti e macrofagi. Supporta e vascolarizza il tessuto epiteliale sovrrastante; è sede di riparazione tissutale.

\subsubsection{Tessuto connettivo fibroso denso}
Ricco di fibre collagene, scarsa sostanza fondamentale, prevalenza di fibrociti. Si distingue in:
\begin{itemize}
  \item \textit{Regolare} (fasci paralleli): caratteristico di tendini e legamenti, massimamente resistente alla trazione unidirezionale;
  \item \textit{Irregolare} (fasci incrociati): presente nel derma e nelle capsule degli organi, resiste a forze multidirezionali.
\end{itemize}

\subsubsection{Tessuto adiposo}
\begin{itemize}
  \item \textbf{Uniloculare} (bianco/giallo): adipociti con un'unica grande goccia lipidica; presente nell'adulto con funzioni di riserva energetica, isolamento termico e sostegno meccanico;
  \item \textbf{Multiloculare} (bruno): adipociti con numerose microgocce e abbondanti mitocondri contenenti \textit{termogenina} (proteina disaccoppiante che genera calore invece di \atp); tipico dei mammiferi ibernanti e dei neonati.
\end{itemize}

% ============================================================
\section{Tessuto cartilagineo}
% ============================================================

La cartilagine è un tessuto di sostegno \textbf{avascolare}, caratterizzato da:
\begin{itemize}
  \item Condrociti nelle lacune della matrice;
  \item Abbondante matrice extracellulare ricca di proteoglicani e fibre.
\end{itemize}

\begin{tcolorbox}[notebox, title={Meccanismo di resistenza alla compressione}]
  I proteoglicani (ricchi di cariche negative) attraggono acqua per effetto osmotico. Sotto carico, l'acqua fuoriesce dalle lacune, avvicinando le cariche negative tra loro e aumentando la resistenza alla compressione. Rimosso il carico, l'acqua rientra e la cartilagine torna alla forma originale.
\end{tcolorbox}

\subsection{Tipi di cartilagine}

\begin{center}
\begin{tabular}{llll}
\toprule
\textbf{Tipo} & \textbf{Fibre prevalenti} & \textbf{Caratteristiche} & \textbf{Sedi tipiche} \\
\midrule
Ialina & Collagene (tipo II) & Liscia, traslucida & Cartilagini articolari, costali, trachea \\
Elastica & Elastiche + collagene & Flessibile, giallognola & Padiglione auricolare, epiglottide \\
Fibrosa & Collagene (tipo I, abbondante) & Rigida, biancastra & Dischi intervertebrali, menischi \\
\bottomrule
\end{tabular}
\end{center}

% ============================================================
\section{Tessuto osseo}
% ============================================================

Il tessuto osseo è un tessuto connettivo specializzato con matrice mineralizzata che conferisce durezza e resistenza. Le sue funzioni principali sono:
\begin{itemize}
  \item Sostegno strutturale e protezione degli organi interni;
  \item Sistema di leve per il movimento;
  \item Riserva di calcio e altri ioni;
  \item Sede del tessuto emopoietico.
\end{itemize}

\subsection{Cellule del tessuto osseo}

\begin{itemize}
  \item \textbf{Osteoblasti:} sintetizzano la matrice organica (osteoide) e promuovono la mineralizzazione; originano dagli osteoprogenitori;
  \item \textbf{Osteociti:} osteoblasti intrappolati nella matrice mineralizzata; mantengono la matrice ossea e partecipano alla riparazione tissutale tramite canalicoli e gap junctions;
  \item \textbf{Osteoclasti:} cellule giganti plurinucleate derivate dai macrofagi; riassorbono la matrice ossea (demineralizzazione + digestione enzimatica lisosomiale).
\end{itemize}

\subsection{Classificazione del tessuto osseo}

\begin{itemize}
  \item \textbf{Osso compatto (lamellare):} formato da osteoni (sistemi di Havers); lamelle ossee concentriche attorno al canale di Havers (vasi nutritizi); resistente alla pressione longitudinale; presente nella diafisi delle ossa lunghe;
  \item \textbf{Osso spugnoso (trabecolare):} lamelle disposte a trabecole con cavità occupate dal midollo osseo; resiste a sollecitazioni multidirezionali; presente nelle epifisi delle ossa lunghe e nelle ossa brevi.
\end{itemize}

\subsection{Ossificazione}

\begin{itemize}
  \item \textbf{Ossificazione diretta} (intramembranosa): il tessuto osseo si forma direttamente da cellule mesenchimali che si differenziano in osteoblasti; caratteristica delle ossa piatte del cranio;
  \item \textbf{Ossificazione indiretta} (endocondrale): un modello cartilagineo viene progressivamente sostituito da tessuto osseo; caratteristica delle ossa lunghe degli arti.
\end{itemize}

% ============================================================
\section{Tessuto emolinfopoietico}
% ============================================================

\subsection{Il sangue}

Il \textbf{sangue} è un tessuto connettivo liquido (pH 7,3--7,4; densità 1,055--1,065\,g/cm$^3$) che costituisce circa il 7\% del peso corporeo (5--6\,L nell'adulto). È composto da:
\begin{itemize}
  \item \textbf{Plasma} (55\%): soluzione acquosa contenente albumina, globuline, fibrinogeno, elettroliti, glucosio, amminoacidi e ormoni;
  \item \textbf{Elementi figurati} (45\%): globuli rossi (99\%), globuli bianchi e piastrine ($<$1\%).
\end{itemize}

\subsubsection{Globuli rossi (eritrociti)}
Cellule anucleate, a forma di disco biconcavo ($\phi \sim 7$\,$\mu$m), contenenti \textbf{emoglobina} (proteina a 4 subunità, ognuna con un gruppo eme che lega un Fe$^{2+}$). Trasportano \Oc{} (come ossiemoglobina) e \CO{}. Vita media $\sim$120 giorni; prodotti per eritropoiesi nel midollo osseo.

\subsubsection{Gruppi sanguigni (sistema ABO)}

\begin{center}
\begin{tabular}{lll}
\toprule
\textbf{Gruppo} & \textbf{Antigeni (emazie)} & \textbf{Anticorpi (plasma)} \\
\midrule
A & Antigene A & Anti-B \\
B & Antigene B & Anti-A \\
AB & A e B & Nessuno \\
0 & Nessuno & Anti-A e anti-B \\
\bottomrule
\end{tabular}
\end{center}

Oltre al sistema ABO, il \textbf{fattore Rh} è determinato dalla presenza (Rh$^+$) o assenza (Rh$^-$) di un antigene specifico sulla superficie dei globuli rossi.

\subsubsection{Piastrine (trombociti)}
Frammenti citoplasmatici anucleati derivati dai megacariociti; fondamentali per l'\textbf{emostasi} (blocco del sanguinamento in caso di lesione vascolare), che avviene in tre fasi:
\begin{enumerate}
  \item \textit{Fase vascolare:} vasocostrizione riflessa del vaso leso;
  \item \textit{Fase piastrinica:} adesione e aggregazione delle piastrine formando un tappo emostatico primario;
  \item \textit{Fase coagulativa:} conversione del fibrinogeno in fibrina (tramite la cascata coagulativa), formando il coagulo definitivo.
\end{enumerate}

\subsubsection{Globuli bianchi (leucociti)}
Cellule nucleate con funzioni di difesa immunitaria. Si classificano in:
\begin{itemize}
  \item \textbf{Granulociti:}
    \begin{itemize}
      \item \textit{Neutrofili} (50--70\%): prima linea di difesa antimicrobica (fagocitosi);
      \item \textit{Eosinofili} (2--4\%): implicati nelle risposte allergiche e antiparassitarie;
      \item \textit{Basofili} (0,5--1\%): rilasciano istamina ed eparina; coinvolti nello shock anafilattico;
    \end{itemize}
  \item \textbf{Non granulari:}
    \begin{itemize}
      \item \textit{Monociti} (2--8\%): si differenziano in macrofagi nei tessuti; attività fagocitaria;
      \item \textit{Linfociti T:} immunità cellulo-mediata (difesa contro microbi intracellulari);
      \item \textit{Linfociti B:} immunità umorale (produzione di anticorpi contro microbi extracellulari);
      \item \textit{Linfociti NK:} sorveglianza immunitaria contro cellule tumorali e viralmente infette.
    \end{itemize}
\end{itemize}

% ============================================================
\section{Tessuto muscolare}
% ============================================================

Il \textbf{tessuto muscolare} è specializzato nella contrazione, che permette il movimento del corpo, la circolazione del sangue e la peristalsi viscerale. È caratterizzato da:
\begin{itemize}
  \item \textbf{Eccitabilità:} risposta agli stimoli nervosi e ormonali;
  \item \textbf{Contrattilità:} capacità di accorciarsi generando forza;
  \item \textbf{Estendibilità ed elasticità.}
\end{itemize}

\begin{center}
\begin{tabular}{llll}
\toprule
\textbf{Tipo} & \textbf{Struttura} & \textbf{Controllo} & \textbf{Sede} \\
\midrule
Scheletrico striato & Multinucleato, striato & Volontario & Muscoli scheletrici \\
Cardiaco striato & Mononucleato, ramificato & Involontario & Cuore \\
Liscio viscerale & Mononucleato, fusiforme & Involontario & Organi cavi, vasi \\
\bottomrule
\end{tabular}
\end{center}

\subsection{Struttura del muscolo scheletrico}

La gerarchia strutturale è:
\[
\text{Fibre muscolari} \;\subset\; \text{Fascicoli (endomisio/perimisio)} \;\subset\; \text{Muscolo (epimisio)}
\]

L'\textbf{unità fondamentale} del tessuto muscolare è il \textbf{sarcomero}: la regione compresa tra due linee Z successive. È composto da:
\begin{itemize}
  \item \textbf{Filamenti sottili} (actina, tropomiosina, troponina): formano le semibande I chiare;
  \item \textbf{Filamenti spessi} (miosina): formano la banda A scura centrale.
\end{itemize}

\subsection{Meccanismo della contrazione muscolare}

La contrazione muscolare segue il \textbf{modello dello scivolamento dei filamenti}:

\begin{enumerate}
  \item Un potenziale d'azione percorre il sarcolemma e si propaga nei tubuli T;
  \item La depolarizzazione induce il rilascio di \Ca{} dal reticolo sarcoplasmatico;
  \item Lo ione \Ca{} si lega alla troponina, inducendo lo spostamento della tropomiosina e l'esposizione dei siti di legame per la miosina sull'actina;
  \item Le teste della miosina si legano all'actina e compiono il \textbf{colpo di forza} (power stroke), facendo scorrere i filamenti sottili verso il centro del sarcomero;
  \item L'idrolisi dell'ATP distacca la testa della miosina dall'actina, permettendo un nuovo ciclo.
\end{enumerate}

\begin{tcolorbox}[defbox, title={Ciclo dei ponti actomiosinici}]
  \begin{enumerate}
    \item $\text{M} + \text{ATP} \to \text{M}\cdot\text{ATP}$ (distacco se legata ad actina)
    \item $\text{M}\cdot\text{ATP} \to \text{M}'\cdot\text{ADP}\cdot P_i$ (idrolisi con sollevamento della testa)
    \item $\text{A} + \text{M}'\cdot\text{ADP}\cdot P_i \to \text{A}\cdot\text{M}'\cdot\text{ADP}\cdot P_i$ (legame forte con actina e rilascio di $P_i$)
    \item Power stroke: $\text{A}\cdot\text{M}'\cdot\text{ADP} \to \text{A}\cdot\text{M}'$ (rilascio di ADP e scivolamento filamenti)
    \item Nuovo legame di ATP $\to$ distacco actina--miosina
  \end{enumerate}
  In assenza di ATP, il ciclo si arresta nel legame forte actomiosinico (\textbf{rigor mortis}).
\end{tcolorbox}

% ============================================================
\section{Tessuto nervoso}
% ============================================================

Il \textbf{tessuto nervoso} riceve stimoli dall'ambiente interno ed esterno, li elabora e trasmette risposte agli organi effettori. È composto da:
\begin{itemize}
  \item \textbf{Neuroni:} cellule eccitabili specializzate nella trasmissione dei segnali;
  \item \textbf{Neuroglia:} cellule di sostegno, nutrizione e mielinizzazione.
\end{itemize}

\subsection{Il neurone}

Il \textbf{neurone} è l'unità funzionale del sistema nervoso e comprende:
\begin{enumerate}
  \item \textbf{Corpo cellulare (soma o pirenoforo):} contiene il nucleo e la maggior parte degli organuli; sede del metabolismo cellulare;
  \item \textbf{Dendriti:} prolungamenti ramificati che ricevono segnali afferenti verso il soma; ospitano spine dendritiche (siti sinaptici);
  \item \textbf{Assone:} unico prolungamento efferente che conduce il potenziale d'azione lontano dal corpo cellulare; può essere lungo fino a oltre un metro.
\end{enumerate}

\subsection{Potenziale d'azione}

Il \textbf{potenziale d'azione} è il segnale elettrico con cui il neurone trasmette l'informazione. Nasce quando uno stimolo depolarizzante raggiunge la soglia di eccitazione al cono di emergenza dell'assone.

Le fasi del potenziale d'azione sono:
\begin{enumerate}
  \item \textbf{Riposo} ($-70$\,mV): canali del K$^+$ aperti; canali del Na$^+$ chiusi (inattivazione);
  \item \textbf{Depolarizzazione:} apertura dei canali voltaggio-dipendenti del Na$^+$; il potenziale sale fino a $+20$\,mV (\textit{overshoot});
  \item \textbf{Ripolarizzazione:} chiusura dei canali del Na$^+$, apertura dei canali del K$^+$; il potenziale torna ai valori negativi;
  \item \textbf{Iperpolarizzazione:} il flusso tardivo di K$^+$ spinge il potenziale sotto il valore di riposo; ingresso di Cl$^-$ per correggere la polarità.
\end{enumerate}

La conduzione è più rapida nelle \textbf{fibre mieliniche}, dove l'impulso salta da un nodo di Ranvier all'altro (\textit{conduzione saltatoria}), rispetto alle fibre amieliniche dove la depolarizzazione si propaga lungo tutta la membrana.

\subsection{Sinapsi}

La \textbf{sinapsi} è la struttura specializzata per la trasmissione del segnale nervoso tra due neuroni o tra un neurone e una cellula effettrice (muscolare, ghiandolare). Si distinguono:

\begin{itemize}
  \item \textbf{Sinapsi elettriche:} gap junctions tra le membrane pre e postsinaptica; trasmissione rapida e bidirezionale;
  \item \textbf{Sinapsi chimiche:} il potenziale d'azione arriva alla terminazione presinaptica, provoca l'ingresso di \Ca{} e l'esocitosi dei neurotrasmettitori nella fessura sinaptica; il neurotrasmettitore si lega ai recettori postsinaptici, generando un nuovo potenziale post-sinaptico.
\end{itemize}

\subsection{Neuroglia}

Nel \textbf{sistema nervoso centrale}:
\begin{itemize}
  \item \textbf{Astrociti:} sostengono i neuroni, mantengono la barriera emato-encefalica, riparano i danni;
  \item \textbf{Oligodendrociti:} mielinizzano gli assoni del SNC (ogni cellula mielinizza più assoni);
  \item \textbf{Cellule ependimali:} rivestono le cavità del sistema nervoso, producono liquido cefalorachidiano;
  \item \textbf{Microglia:} macrofagi residenti del SNC; attivati da infiammazione o danno.
\end{itemize}

Nel \textbf{sistema nervoso periferico}:
\begin{itemize}
  \item \textbf{Cellule di Schwann:} mielinizzano i singoli assoni periferici; permettono la rigenerazione nervosa;
  \item \textbf{Cellule satelliti:} circondano i corpi cellulari nei gangli; regolano l'omeostasi dei neurotrasmettitori.
\end{itemize}

% ============================================================
\backmatter
% ============================================================

\cleardoublepage
\phantomsection
\addcontentsline{toc}{chapter}{Indice analitico}

\end{document}